\documentclass[a4paper]{article}
\usepackage[fontset=fandol]{ctex}
\usepackage[margin=1.65cm]{geometry}
\usepackage{amsmath,amssymb}
\usepackage{graphicx}
\usepackage{hyperref}
\usepackage{longtable,booktabs}
\usepackage{xcolor}
\usepackage{enumitem}
\usepackage{listings}
\lstset{basicstyle=\ttfamily\small,breaklines=true,columns=fullflexible,keepspaces=true}
\hypersetup{colorlinks=true,linkcolor=blue,urlcolor=blue}
\setlist[itemize]{leftmargin=1.25em,itemsep=0.18em}
\setlength{\parskip}{0.45em}
\setlength{\parindent}{2em}
\sloppy
\title{第 11 讲：Scaling laws 进阶与实践}
\author{CS336 Spring 2026 public videos and official course materials}
\date{}
\begin{document}
\maketitle
\tableofcontents
\newpage
\section{本章导读}
本章讨论 \textbf{Scaling laws 进阶与实践}。CS336 的第一层目标不是把现成大模型当作黑箱使用，而是把 language model 从输入文本、训练目标、模型结构、数据、系统和评估这些部件重新拆开。只有拆开之后，读者才能判断一个设计选择究竟改变了什么。
设想基础 scaling law 给出一条漂亮曲线，但你的数据有限、训练 schedule 改了、benchmark 也不只看 loss。进阶 scaling 的问题是：哪些外推仍可信，哪些必须重新校准？
本章对应 CS336 Spring 2026 第 11 个公开视频，并结合课程页提供的可解析资料。视频和课程页给出事实边界；正文把这些材料改写为可自学的教材章节。
阅读本章时，先理解问题为什么存在，再看定义、公式和实现形态，随后通过例题和练习检查自己是否能独立复现这条 reasoning chain。
\begin{figure}[h]
\centering
\includegraphics[width=0.58\linewidth]{source_anchor.jpg}
\caption{第 11 讲的课程图像锚点。}
\end{figure}
\subsection{本章知识路线}
\begin{itemize}
\item 本章首先说明 Scaling laws 进阶与实践 为什么是 language modeling pipeline 中不可跳过的一环。
\item 其次定义本章反复使用的术语，并说明这些术语对应的计算对象或系统对象。
\item 随后用公式和伪代码表达主要机制，让概念能够落到可计算的变量上。
\item 最后给出例题、风险讨论和练习，便于读者检查自己是否真正掌握了机制。
\end{itemize}
\subsection{结构图}
下面的图用于帮助读者定位本章的核心结构。先看概念之间的依赖，再回到正文中的公式、伪代码和例题。
\begin{figure}[h]
\centering
\includegraphics[width=0.72\linewidth]{pdf_pages/page-02.jpg}
\caption{官方材料图页：Scaling laws 进阶与实践 的课程页/PDF/script 图像锚点。}
\end{figure}
\begin{figure}[h]
\centering

\includegraphics[width=0.72\linewidth]{pdf_pages/page-01.jpg}
\caption{官方材料图页：Scaling laws 进阶与实践 的课程页/PDF/script 图像锚点。}
\end{figure}
\section{本章主线}
Scaling laws 进阶与实践 的主线是从小规模实验走向大规模决策。读者需要同时追踪参数量、token 数、compute budget、loss 和 benchmark 之间的关系。
\subsection{本章要解决的核心问题}
本讲继续 scaling laws，重点转向实践细节：muP、WSD learning rate、batch/LR scaling、MiniCPM 和 DeepSeek 等公开案例。
读完这一段后，应能用一两句话说明 Scaling laws 进阶与实践 要解决的瓶颈，并指出这个瓶颈发生在数据、模型、优化、系统还是评估层。
\subsection{从机制到资源账本}
进阶 scaling 的难点是如何降低拟合 scaling law 本身的成本。WSD 等 schedule 允许从稳定阶段重启 decay，以便更便宜地估计数据-模型 tradeoff。
随后把机制写成账本：输入规模是什么，主要状态是什么，成本或质量指标怎样随变量改变。这个账本让后面的公式不再只是符号。
\subsection{从课堂讲解到可复现实验}
实践中要同时预测 loss、选择 batch、学习率、模型宽深比和训练 tokens；任何 recipe 改变都需要重新验证外推可靠性。
最后把课堂讲解转成可复现实验：固定协议，改变一个变量，记录中间量，并解释结果为何支持或反驳原来的机制判断。
\section{术语、符号与问题边界}
本章的术语横跨模型、数据和系统。读者需要先知道每个词指向的对象：有些词描述数学目标，有些词描述张量形状，有些词描述硬件瓶颈，还有一些词描述评估协议。术语边界不清，后面的公式和实现就会变成空洞的缩写。
\begin{longtable}{p{0.08\linewidth}p{0.42\linewidth}p{0.42\linewidth}}
\toprule
\textbf{\#} & \textbf{术语} & \textbf{阅读要求}\\
\midrule
1 & forecasting（预测） & 给出定义、一个最小例子、一个关联公式或代码位置，以及一个典型失败情形。\\
2 & overtraining（过训练/超额训练） & 给出定义、一个最小例子、一个关联公式或代码位置，以及一个典型失败情形。\\
3 & data-constrained scaling（数据受限缩放） & 说明它如何改变训练样本或 token 分布，并给出一个会改变结论的数据边界条件。\\
4 & muP & 给出定义、一个最小例子、一个关联公式或代码位置，以及一个典型失败情形。\\
5 & emergence（涌现） & 给出定义、一个最小例子、一个关联公式或代码位置，以及一个典型失败情形。\\
\bottomrule
\end{longtable}
这些术语之间有依赖关系。例如 tokenization 改变 sequence length，sequence length 又改变 attention FLOPs、KV cache 和 evaluation token budget；parallelism 改变显存占用，也改变通信瓶颈和故障恢复方式。
\subsection{术语之间的关系}
在本章中，forecasting（预测） 给出问题入口，overtraining（过训练/超额训练） 描述主要机制，data-constrained scaling（数据受限缩放） 则把机制连接到可测量的结果。读教材时应当沿着这个链条追问：输入是什么，中间状态怎样改变，输出指标为什么随之变化。
课程材料还会反复出现 scaling、batch、optimal、scale、learning、chinchilla 等词。它们不是一组同义标签，而是提示本章的不同层次：有的指对象，有的指操作，有的指约束或指标。
因此，术语表不是背诵清单，而是后文阅读的索引。遇到一个术语时，先定位它属于 representation、optimization、systems、data 还是 evaluation，再判断它改变哪一项账本。
\section{更复杂的 scaling 问题}
现在进入 更复杂的 scaling 问题。这一节把模型大小、训练 token、compute budget 和 loss 放在同一坐标系里，讨论何时可以从小实验推断大训练。
\subsection{为什么要讨论 更复杂的 scaling 问题}
视频开场称本讲是 scaling journey 的继续，覆盖更高级细节和实践中扩展模型的重要问题。
实践中不只问 loss 怎么随 \textbackslash{}(N,D\textbackslash{}) 变，还要问数据受限、训练长度、architecture 改变、optimizer 改变时预测是否仍成立。
scaling law 的价值在于减少大 run 风险，而不是替代工程监控。
更复杂的 scaling 问题 把 Scaling laws 进阶与实践 连接到大规模决策。接下来要看的不是曲线是否平滑，而是实验点、预算变量和外推假设是否一致。
\subsection{更复杂的 scaling 问题 的机制}
视频开场称本讲是 scaling journey 的继续，覆盖更高级细节和实践中扩展模型的重要问题。 在 更复杂的 scaling 问题 中，核心对象是模型规模、token 数和训练预算：小实验只有在预算、数据和优化条件清楚时，才可能支持大规模外推。
实践中不只问 loss 怎么随 \textbackslash{}(N,D\textbackslash{}) 变，还要问数据受限、训练长度、architecture 改变、optimizer 改变时预测是否仍成立。 因而同一个方法在不同规模下可能改变不同的变量，尤其是参数量、训练 token 数、compute budget、loss 外推误差和 benchmark 转化关系。比较方法时，必须同时给出问题规模和实验条件。
scaling law 的价值在于减少大 run 风险，而不是替代工程监控。 可检验性来自多个实验点、残差分析和外推误差估计；单条曲线若没有不确定性，就不能支撑大规模决策。
因此，更复杂的 scaling 问题 的学习目标是建立一套外推判断：读者应能说明小实验为何能或不能指导 Scaling laws 进阶与实践 中的大训练配置。
\subsection{更复杂的 scaling 问题 的数学表达}
\[L(C)=L_\infty + aC^{-\gamma}\]
符号说明：\(C\) 是总 compute；这个一维形式常用于简化预测，但会隐藏 \(N,D\) tradeoff。
这个关系式给出 更复杂的 scaling 问题 的缩放假设。它把参数量、训练 token 数和 compute budget 放到同一张账本中，重点是判断小实验能否可靠外推到大训练。
在 scaling 实验中，validation loss、benchmark transfer 和训练吞吐都只是不同侧面的 proxy。外推时要说明哪个 proxy 被拟合，哪个 proxy 只是辅助诊断。
\subsection{从 更复杂的 scaling 问题 到程序}
把 更复杂的 scaling 问题 写成程序后，抽象机制会变成输入、状态和输出之间的变换。下面的伪代码保留核心计算顺序。
\begin{lstlisting}[language=Python]
fit_compute_curve(runs)
check_breakpoints_for_new_regime()
forecast_only_with_error_bars()
\end{lstlisting}
阅读这段流程时，重点看小规模实验怎样产生可外推曲线。参数量、训练 token 数、compute budget、loss 外推误差和 benchmark 转化关系 需要在每个实验点上同时记录，缺少其中一项都会改变 scaling 解释。
\subsection{边界条件与常见误解}
\begin{itemize}
\item 单变量 compute curve 可能掩盖模型大小和数据量的不同失败模式。
\item architecture 改变后，把旧曲线直接外推很危险。
\item 小规模 scaling 曲线不一定覆盖大训练的优化不稳定、数据瓶颈和系统瓶颈。
\item compute-optimal 结论依赖目标指标；以 loss 最优并不自动意味着 benchmark、latency 或成本最优。
\end{itemize}
\subsection{与本章其他概念的关系}
更复杂的 scaling 问题 连接本章的小实验和大规模决策：术语表标出预算变量，公式给出外推结构，伪代码组织实验点，例题检查 compute-optimal 判断是否稳健。
\subsection{怎样检验 更复杂的 scaling 问题}
检验 更复杂的 scaling 问题 时，先把抽象说法落到模型规模、token 数和训练预算上。与 forecasting、compute、regime、loss、practice 相关的对象必须能被构造、打印、计数或评分；否则实验只能得到描述性观察，不能得到可复现结论。
一个合适的小实验应当保留本节机制，同时让 参数量、训练 token 数、compute budget、loss 外推误差和 benchmark 转化关系 中至少一个量发生可解释变化。输入规模不必逼真；关键是读者能手算或打印关键中间量，并说明哪个变量触发了结果变化。
随后把公式说明转成可记录的数量关系。符号说明：\textbackslash{}(C\textbackslash{}) 是总 compute；这个一维形式常用于简化预测，但会隐藏 \textbackslash{}(N,D\textbackslash{}) tradeoff。 如果数量只能间接观测，就必须说明使用了什么 proxy；如果使用 benchmark 或 reward，也要说明协议和随机性。
scaling 消融通常一次只改变参数量、token 数、数据 pipeline 或优化 schedule。若多个轴同时移动，曲线拟合会混合不同机制。
最后检查失败条件。工程判断往往在反例中变清楚：单变量 compute curve 可能掩盖模型大小和数据量的不同失败模式；architecture 改变后，把旧曲线直接外推很危险。复习时可以把这些限制改写成反例测试，观察它们怎样改变 参数量、训练 token 数、compute budget、loss 外推误差和 benchmark 转化关系。
把实验结论放回 Scaling laws 进阶与实践 时，要说明 更复杂的 scaling 问题 改变的是预算轴、拟合曲线、外推误差还是目标指标。能拆开这些因素，才说明读者真正理解 scaling。
\subsection{概念辨析：更复杂的 scaling 问题}
forecasting、compute、regime、loss、practice 常一起出现在 scaling 讨论中，但它们分别指预算轴、拟合对象、外推假设或风险来源。读者应先确定曲线中的自变量和因变量。
区分这些词时，可以先问它们是否改变 参数量、训练 token 数、compute budget、loss 外推误差和 benchmark 转化关系。一个术语如果不能对应到变量、状态、代码路径或评估协议，就还没有形成可操作含义。
围绕 更复杂的 scaling 问题 复习时，可以画三到五个小实验点：记录参数量、token 数、预算和 loss，再检查外推是否合理。
\subsection{更复杂的 scaling 问题 的小结}
\begin{itemize}
\item 讨论 更复杂的 scaling 问题 时，应同时报告问题规模、实验设置和主要观测变量，尤其是 参数量、训练 token 数、compute budget、loss 外推误差和 benchmark 转化关系。
\item 如果方法声称可以外推，要说明实验点覆盖范围、拟合形式和数据管线是否一致；曲线本身不是自然定律。
\item scaling 结论需要把 loss、benchmark 和成本之间的转换关系说清楚，否则 compute-optimal 只是一条局部经验曲线。
\end{itemize}
\section{数据受限与重复数据}
现在进入 数据受限与重复数据。这一节关心训练样本如何形成、被筛选并进入混合分布；它直接决定模型看到什么，也决定 evaluation 是否可能被污染。
\subsection{为什么要讨论 数据受限与重复数据}
当高质量 token 有限时，简单增加训练 tokens 可能意味着重复数据或低质量数据。
重复训练同一数据会改变 scaling 行为：短期可能降低 loss，长期可能过拟合或降低泛化。
这为后续 data lectures 铺垫：scaling law 必须和数据来源、过滤、去重、混合策略一起讨论。
数据受限与重复数据 把 Scaling laws 进阶与实践 连接到训练分布。接下来要看的不是语料名称，而是样本来源、处理顺序和版本记录怎样改变模型可学习的证据。
\subsection{数据受限与重复数据 的机制}
当高质量 token 有限时，简单增加训练 tokens 可能意味着重复数据或低质量数据。 在 数据受限与重复数据 中，核心对象是样本、来源和分布状态：数据管线每改变一次样本集合，模型实际学习的分布也随之改变。
重复训练同一数据会改变 scaling 行为：短期可能降低 loss，长期可能过拟合或降低泛化。 因而同一个方法在不同规模下可能改变不同的变量，尤其是样本数量、去重率、过滤阈值、数据混合比例和 contamination 风险。比较方法时，必须同时给出问题规模和实验条件。
这为后续 data lectures 铺垫：scaling law 必须和数据来源、过滤、去重、混合策略一起讨论。 可检验性来自 provenance 记录、过滤前后计数、抽样审计和 held-out contamination 检查；没有这些记录，数据结论无法复现。
因此，数据受限与重复数据 的学习目标是建立一条数据谱系：读者应能从原始来源追到训练 token，并说明每个处理步骤改变了什么分布。
\subsection{数据受限与重复数据 的数学表达}
\[D_{\text{effective}}\le D_{\text{raw}}\quad\text{when duplicates and low-quality tokens are present}\]
符号说明：\(D_{\text{raw}}\) 是原始 token 数，\(D_{\text{effective}}\) 是对泛化真正有贡献的有效 token 数。
这个关系式刻画 数据受限与重复数据 对数据分布的影响。读者应关心哪些变量来自采集过程，哪些变量来自过滤或混合策略，哪些变量只能通过抽样审计估计。
在数据实验中，去重率、过滤通过率、语言比例和抽样人工判读都是 proxy。它们不能单独代表数据质量，必须和来源记录、版本号以及污染检查一起解释。
\subsection{从 数据受限与重复数据 到程序}
把 数据受限与重复数据 写成程序时，要明确每一步怎样改写样本集合。下面的伪代码强调输入、过滤、统计和输出之间的关系。
\begin{lstlisting}[language=Python]
effective_tokens = estimate_after_dedup_and_filter(raw_tokens)
refit_scaling_law_if_data_pipeline_changes()
\end{lstlisting}
阅读这段流程时，重点看 数据受限与重复数据 怎样改变样本集合。每个过滤、去重或混合步骤都应留下计数和抽样检查，使 样本数量、去重率、过滤阈值、数据混合比例和 contamination 风险 能够被复核。
\subsection{边界条件与常见误解}
\begin{itemize}
\item 有效 token 数不可直接从文件大小读取。
\item 低质量数据可能增加 token count 但恶化 downstream behavior。
\item 小样本抽样只能发现部分数据问题，不能证明语料整体无污染或无偏差。
\item 数据版本、许可状态和处理脚本如果缺失，后续 evaluation 与 scaling 结论都会失去可复现基础。
\end{itemize}
\subsection{与本章其他概念的关系}
数据受限与重复数据 连接本章的数据来源和模型行为：定义说明样本如何进入语料，公式刻画分布变化，伪代码展示处理顺序，例题检验数据决策怎样影响训练结论。
\subsection{怎样检验 数据受限与重复数据}
检验 数据受限与重复数据 时，先把抽象说法落到样本、来源和分布状态上。与 data、duplicates、effective tokens、overtraining、quality 相关的对象必须能被构造、打印、计数或评分；否则实验只能得到描述性观察，不能得到可复现结论。
一个合适的小实验应当保留本节机制，同时让 样本数量、去重率、过滤阈值、数据混合比例和 contamination 风险 中至少一个量发生可解释变化。输入规模不必逼真；关键是读者能手算或打印关键中间量，并说明哪个变量触发了结果变化。
随后把公式说明转成可记录的数量关系。符号说明：\textbackslash{}(D\_\{ ext\{raw\}\}\textbackslash{}) 是原始 token 数，\textbackslash{}(D\_\{ ext\{effective\}\}\textbackslash{}) 是对泛化真正有贡献的有效 token 数。 如果数量只能间接观测，就必须说明使用了什么 proxy；如果使用 benchmark 或 reward，也要说明协议和随机性。
数据消融通常一次只改变过滤阈值、去重策略、混合权重或数据版本。若同时更换 tokenizer、语料和评估集，就无法知道改进来自哪一步。
最后检查失败条件。工程判断往往在反例中变清楚：有效 token 数不可直接从文件大小读取；低质量数据可能增加 token count 但恶化 downstream behavior。复习时可以把这些限制改写成反例测试，观察它们怎样改变 样本数量、去重率、过滤阈值、数据混合比例和 contamination 风险。
把实验结论放回 Scaling laws 进阶与实践 时，要说明 数据受限与重复数据 改变的是样本分布、数据质量、污染风险还是可复现性。能追到这些来源，才说明读者真正理解数据机制。
\subsection{概念辨析：数据受限与重复数据}
data、duplicates、effective tokens、overtraining、quality 常一起出现在数据管线中，但它们分别指来源、处理动作、版本记录或风险标记。读者应先判断它们改变的是样本集合还是审计证据。
其中，data（数据）是模型能力、偏差和法律风险的主要来源。
区分这些词时，可以先问它们是否改变 样本数量、去重率、过滤阈值、数据混合比例和 contamination 风险。一个术语如果不能对应到变量、状态、代码路径或评估协议，就还没有形成可操作含义。
围绕 数据受限与重复数据 复习时，可以构造一小批文档：保留来源和版本，运行一个处理步骤，再比较样本计数和抽样质量。
\subsection{数据受限与重复数据 的小结}
\begin{itemize}
\item 讨论 数据受限与重复数据 时，应同时报告问题规模、实验设置和主要观测变量，尤其是 样本数量、去重率、过滤阈值、数据混合比例和 contamination 风险。
\item 如果方法声称数据更好，要说明过滤、去重、混合或合成步骤改变了哪一部分分布，并检查 benchmark contamination。
\item 数据结论需要抽样审计和版本记录；没有 provenance 的语料变化不能支持可引用结论。
\end{itemize}
\section{从 loss 到能力指标}
现在进入 从 loss 到能力指标。这一节把模型大小、训练 token、compute budget 和 loss 放在同一坐标系里，讨论何时可以从小实验推断大训练。
\subsection{为什么要讨论 从 loss 到能力指标}
loss 平滑、可测、适合拟合，但用户关心的是 reasoning、coding、instruction following 等能力。
能力指标可能出现阈值效应或评估噪声，使其比 loss 更难预测。
进阶 scaling 分析要把 validation loss、benchmark score、推理成本和安全约束一起纳入。
从 loss 到能力指标 把 Scaling laws 进阶与实践 连接到大规模决策。接下来要看的不是曲线是否平滑，而是实验点、预算变量和外推假设是否一致。
\subsection{从 loss 到能力指标 的机制}
loss 平滑、可测、适合拟合，但用户关心的是 reasoning、coding、instruction following 等能力。 在 从 loss 到能力指标 中，核心对象是模型规模、token 数和训练预算：小实验只有在预算、数据和优化条件清楚时，才可能支持大规模外推。
能力指标可能出现阈值效应或评估噪声，使其比 loss 更难预测。 因而同一个方法在不同规模下可能改变不同的变量，尤其是参数量、训练 token 数、compute budget、loss 外推误差和 benchmark 转化关系。比较方法时，必须同时给出问题规模和实验条件。
进阶 scaling 分析要把 validation loss、benchmark score、推理成本和安全约束一起纳入。 可检验性来自多个实验点、残差分析和外推误差估计；单条曲线若没有不确定性，就不能支撑大规模决策。
因此，从 loss 到能力指标 的学习目标是建立一套外推判断：读者应能说明小实验为何能或不能指导 Scaling laws 进阶与实践 中的大训练配置。
\subsection{从 loss 到能力指标 的数学表达}
\[\mathbb{E}[\text{score}]\approx h(L,\text{prompt},\text{eval distribution})\]
符号说明：\(h\) 表示从 loss 到 benchmark score 的经验映射；它依赖评估任务和 prompting。
这个关系式给出 从 loss 到能力指标 的缩放假设。它把参数量、训练 token 数和 compute budget 放到同一张账本中，重点是判断小实验能否可靠外推到大训练。
在 scaling 实验中，validation loss、benchmark transfer 和训练吞吐都只是不同侧面的 proxy。外推时要说明哪个 proxy 被拟合，哪个 proxy 只是辅助诊断。
\subsection{从 从 loss 到能力指标 到程序}
把 从 loss 到能力指标 写成程序后，抽象机制会变成输入、状态和输出之间的变换。下面的伪代码保留核心计算顺序。
\begin{lstlisting}[language=Python]
for checkpoint in training_curve:
    measure_loss(checkpoint)
    measure_benchmarks(checkpoint)
    fit_loss_to_score_mapping()
\end{lstlisting}
阅读这段流程时，重点看小规模实验怎样产生可外推曲线。参数量、训练 token 数、compute budget、loss 外推误差和 benchmark 转化关系 需要在每个实验点上同时记录，缺少其中一项都会改变 scaling 解释。
\subsection{边界条件与常见误解}
\begin{itemize}
\item benchmark score 的方差和 contamination 会扭曲趋势。
\item 能力不是只由 pretraining loss 决定，post-training 也会改变可见行为。
\item 小规模 scaling 曲线不一定覆盖大训练的优化不稳定、数据瓶颈和系统瓶颈。
\item compute-optimal 结论依赖目标指标；以 loss 最优并不自动意味着 benchmark、latency 或成本最优。
\end{itemize}
\subsection{与本章其他概念的关系}
从 loss 到能力指标 连接本章的小实验和大规模决策：术语表标出预算变量，公式给出外推结构，伪代码组织实验点，例题检查 compute-optimal 判断是否稳健。
\subsection{怎样检验 从 loss 到能力指标}
检验 从 loss 到能力指标 时，先把抽象说法落到模型规模、token 数和训练预算上。与 benchmark、emergence、score、loss、post-training 相关的对象必须能被构造、打印、计数或评分；否则实验只能得到描述性观察，不能得到可复现结论。
一个合适的小实验应当保留本节机制，同时让 参数量、训练 token 数、compute budget、loss 外推误差和 benchmark 转化关系 中至少一个量发生可解释变化。输入规模不必逼真；关键是读者能手算或打印关键中间量，并说明哪个变量触发了结果变化。
随后把公式说明转成可记录的数量关系。符号说明：\textbackslash{}(h\textbackslash{}) 表示从 loss 到 benchmark score 的经验映射；它依赖评估任务和 prompting。 如果数量只能间接观测，就必须说明使用了什么 proxy；如果使用 benchmark 或 reward，也要说明协议和随机性。
scaling 消融通常一次只改变参数量、token 数、数据 pipeline 或优化 schedule。若多个轴同时移动，曲线拟合会混合不同机制。
最后检查失败条件。工程判断往往在反例中变清楚：benchmark score 的方差和 contamination 会扭曲趋势；能力不是只由 pretraining loss 决定，post-training 也会改变可见行为。复习时可以把这些限制改写成反例测试，观察它们怎样改变 参数量、训练 token 数、compute budget、loss 外推误差和 benchmark 转化关系。
把实验结论放回 Scaling laws 进阶与实践 时，要说明 从 loss 到能力指标 改变的是预算轴、拟合曲线、外推误差还是目标指标。能拆开这些因素，才说明读者真正理解 scaling。
\subsection{概念辨析：从 loss 到能力指标}
benchmark、emergence、score、loss、post-training 常一起出现在 scaling 讨论中，但它们分别指预算轴、拟合对象、外推假设或风险来源。读者应先确定曲线中的自变量和因变量。
区分这些词时，可以先问它们是否改变 参数量、训练 token 数、compute budget、loss 外推误差和 benchmark 转化关系。一个术语如果不能对应到变量、状态、代码路径或评估协议，就还没有形成可操作含义。
围绕 从 loss 到能力指标 复习时，可以画三到五个小实验点：记录参数量、token 数、预算和 loss，再检查外推是否合理。
\subsection{从 loss 到能力指标 的小结}
\begin{itemize}
\item 讨论 从 loss 到能力指标 时，应同时报告问题规模、实验设置和主要观测变量，尤其是 参数量、训练 token 数、compute budget、loss 外推误差和 benchmark 转化关系。
\item 如果方法声称可以外推，要说明实验点覆盖范围、拟合形式和数据管线是否一致；曲线本身不是自然定律。
\item scaling 结论需要把 loss、benchmark 和成本之间的转换关系说清楚，否则 compute-optimal 只是一条局部经验曲线。
\end{itemize}
\section{综合讨论：从局部机制到完整系统}
本章的主题是 Scaling laws 进阶与实践，但它不应被读成几个相邻知识点的拼接。更复杂的 scaling 问题、数据受限与重复数据、从 loss 到能力指标 共同回答同一个问题：语言模型系统中哪些对象可以被定义，哪些成本可以被估算，哪些结论必须通过实验或评估来确认。
这些对象包括 forecasting（预测）、overtraining（过训练/超额训练）、data-constrained scaling（数据受限缩放）、muP、emergence（涌现）。读者可以把它们看成一组接口：有的接口连接文本和 token，有的连接模型结构和张量计算，有的连接数据样本和训练目标，有的连接离线评估和在线服务。接口一旦改变，后面的成本、错误类型和优化空间也会随之改变。
\subsection{更复杂的 scaling 问题 在系统中的位置}
更复杂的 scaling 问题 首先提供一个局部解释：视频开场称本讲是 scaling journey 的继续，覆盖更高级细节和实践中扩展模型的重要问题。 这类解释的价值在于把直觉压缩成可操作对象。一旦对象明确，读者就可以追问它的输入、输出、中间状态和失败条件，而不是只记住方法名。
其次，更复杂的 scaling 问题 会改变至少一个资源或评估变量。它可能改变 sequence length、activation size、memory bandwidth、communication volume、data mixture、reward distribution 或 benchmark protocol。这些变量在小例子里可以手算，在真实训练和服务系统里则需要 profiling、logging 和 held-out evaluation。
最后，更复杂的 scaling 问题 不能脱离边界条件理解。一个关键 caveat 是：单变量 compute curve 可能掩盖模型大小和数据量的不同失败模式。。把 caveat 写出来不是为了削弱结论，而是为了说明结论在哪些输入、硬件、数据和评估条件下才成立。
\subsection{数据受限与重复数据 在系统中的位置}
数据受限与重复数据 首先提供一个局部解释：当高质量 token 有限时，简单增加训练 tokens 可能意味着重复数据或低质量数据。 这类解释的价值在于把直觉压缩成可操作对象。一旦对象明确，读者就可以追问它的输入、输出、中间状态和失败条件，而不是只记住方法名。
其次，数据受限与重复数据 会改变至少一个资源或评估变量。它可能改变 sequence length、activation size、memory bandwidth、communication volume、data mixture、reward distribution 或 benchmark protocol。这些变量在小例子里可以手算，在真实训练和服务系统里则需要 profiling、logging 和 held-out evaluation。
最后，数据受限与重复数据 不能脱离边界条件理解。一个关键 caveat 是：有效 token 数不可直接从文件大小读取。。把 caveat 写出来不是为了削弱结论，而是为了说明结论在哪些输入、硬件、数据和评估条件下才成立。
\subsection{从 loss 到能力指标 在系统中的位置}
从 loss 到能力指标 首先提供一个局部解释：loss 平滑、可测、适合拟合，但用户关心的是 reasoning、coding、instruction following 等能力。 这类解释的价值在于把直觉压缩成可操作对象。一旦对象明确，读者就可以追问它的输入、输出、中间状态和失败条件，而不是只记住方法名。
其次，从 loss 到能力指标 会改变至少一个资源或评估变量。它可能改变 sequence length、activation size、memory bandwidth、communication volume、data mixture、reward distribution 或 benchmark protocol。这些变量在小例子里可以手算，在真实训练和服务系统里则需要 profiling、logging 和 held-out evaluation。
最后，从 loss 到能力指标 不能脱离边界条件理解。一个关键 caveat 是：benchmark score 的方差和 contamination 会扭曲趋势。。把 caveat 写出来不是为了削弱结论，而是为了说明结论在哪些输入、硬件、数据和评估条件下才成立。
\subsection{把本章知识用于读论文和复现实验}
读相关论文时，可以把本章变成一个检查表。第一，论文是否清楚定义了输入规模、模型规模、数据来源和评估协议；第二，论文声称的改进落在哪个变量上；第三，作者是否报告了会让结论失效的设置；第四，实验是否能分离模型结构、数据质量、优化设置和系统实现的影响。
复现实验时，也应避免直接追求大规模。先用最小程序复现一个公式、一个伪代码分支或一个 profile 现象，再扩大输入规模。这样做的原因很简单：如果小规模程序无法解释中间量，大规模训练只会把错误隐藏在日志和随机波动里。
把这些检查合在一起，本章给出的不是一个固定答案，而是一种阅读 CS336 的方法：每个模型机制都要落到变量，每个变量都要落到测量，每个测量都要说明协议，每个协议都要承认边界。
因此，本章中的任何一句结论都可以被改写成一个可引用的完整句式：在给定数据、模型、硬件和评估协议下，某个机制改变了某个变量，并通过某个观测指标表现出来。这样的句子比单独背诵术语更长，但它包含了自学、复现和引用时真正需要的信息。
把本章用于自学时，可以按三遍阅读。第一遍只追踪对象和定义，确认每个术语到底指向文本、token、张量、样本、请求还是评估记录。第二遍追踪公式和伪代码，确认每一步改变了哪些状态。第三遍追踪实验和 caveat，确认哪些结论只在特定规模、数据或硬件条件下成立。
把本章用于复习时，则应反过来操作：先从一个失败案例或异常指标出发，再回到相应机制。loss 曲线异常可能指向数据、优化或模型容量；吞吐异常可能指向 kernel、通信或调度；benchmark 异常可能指向 contamination、prompt protocol 或采样设置。能够从异常反推机制，是掌握本章的实用标准。
\section{例题：把概念变成可计算检查}
\subsection{为什么数据 pipeline 改变后要重拟合 scaling law}
同一模型大小和 token 数，换用更强过滤后的数据。
\begin{itemize}
\item validation loss 曲线可能整体下移，因为有效 token 质量提高。
\item 不同 benchmark 的改善幅度不一定一致。
\item 旧曲线的外推误差不再代表新数据分布。
\end{itemize}
\textbf{结论。} scaling law 不是模型大小的自然属性，而是模型、数据、优化和评估分布的联合经验规律。
这个例题的目的不是替代完整工程实现，而是给出一种最小推理形式：把问题转成变量，写出近似公式或伪代码，再说明近似何时失效。
\section{实验设计：从小例子到可复现结论}
学习 Scaling laws 进阶与实践 不能停在概念层。一个教材化结论至少需要四件事：可构造的输入、可观测的中间量、可比较的指标，以及清楚的失败条件。下面的实验设计不是要求训练大模型，而是要求读者用小规模程序复现本章机制。
\subsection{最小输入}
先构造只触发本章核心机制的小输入。例如 tokenizer 章节可以用几个包含 rare words 和 Unicode 字符的字符串；GPU 章节可以用一个 elementwise kernel 和一个 GEMM；data 章节可以用十几篇近重复文档。小输入的价值在于它让错误可见。
\subsection{资源账本}
对同一输入写下 FLOPs、bytes moved、显存峰值、通信量或样本数量的估算。估算不需要一开始就精确，但必须说明单位和近似假设。随后用 profiler、benchmark、validation loss 或人工检查去验证估算是否同量级。
\subsection{单变量消融}
一次只改变一个因素，例如 vocabulary size、head 数、batch size、filter threshold、reward scale、decoding temperature 或 GPU 数量。若多个变量同时变化，实验只能说明系统变了，不能说明机制是什么。
\subsection{失败条件}
最后要刻意触发失败：极端 context length、错误 dtype、过小 batch、污染 benchmark、低质量数据或不可靠 verifier。失败条件常常比成功曲线更能说明一个方法的边界。
\subsection{记录与报告}
实验记录应能让另一个读者复现同一结论。最少要写清楚输入构造、代码版本、随机种子、硬件、batch shape、数据版本、评价指标和异常处理方式。报告结论时不要只写“更快”或“更好”，而要说明快在哪里、好在哪个指标上，以及是否牺牲了内存、稳定性、数据质量或可解释性。
常见报告错误是把局部现象写成普遍规律。例如一次小模型实验里的 loss 改善，不等于更大模型或不同数据分布上也会改善；一次 kernel benchmark 的吞吐提升，不等于端到端 serving latency 一定降低。教材中的实验报告应始终把结论限定在可复现条件内。
\subsection{从小实验到大结论}
小实验的作用不是替代课程中的完整训练或系统实现，而是建立一条可检查的推理链。先在小输入上确认变量方向，再在中等规模上确认数量级，最后才讨论是否能外推到更大模型、更多 token 或更复杂 workload。缺少中间层级时，大结论通常只是在复述直觉。
因此，实验报告最好同时包含正例和反例：正例说明机制在受控条件下怎样工作，反例说明条件稍变时哪里失效。对 CS336 这样的课程，反例尤其重要，因为 language model 的错误常来自多个层面叠加：数据分布、优化稳定性、硬件瓶颈、评估协议和服务负载都可能同时改变观察结果。
还要记录资料版本与复现边界。公开视频、课程页、配套代码和 PDF 可能在不同时间更新；复现实验应写明使用的材料版本，并说明哪些结论来自课程事实，哪些是为了帮助理解而加入的解释性推导。这样才能把学习笔记变成可检查的技术记录。
若复现实验与课程结论不一致，首先检查输入、版本和评估协议，而不是立即修改模型假设。
\subsection{把本章重新读成一条证据链}
完成小实验之后，可以把 Scaling laws 进阶与实践 重新读成一条证据链：先用“更复杂的 scaling 问题”确定问题入口，再用“数据受限与重复数据”解释主要机制，最后用“从 loss 到能力指标”检查边界或外推条件。这种读法比按页背诵更接近教材训练，因为它要求每个结论都能说明前提、变量和证据。
本章反复使用的术语包括 forecasting（预测）、overtraining（过训练/超额训练）、data-constrained scaling（数据受限缩放）、muP、emergence（涌现）。复习时不要只写定义，而要为每个术语补上一句操作性说明：它在课程代码、公式、数据记录或评估协议中对应什么对象；如果这个对象被删除、替换或缩放，哪一个观测量会变化。
最后，用一个反例结束复习。反例可以是资源估算不准、数据污染、reward 失真、benchmark 解析失败、长上下文显存爆炸或 scaling 外推失效。能解释反例的章节，才真正具备自学和引用价值；不能解释反例的章节，只是把课程材料换了一种排版。
引用本章结论时，还应保留来源边界：哪些说法直接来自视频、课程页、slides 或代码，哪些是为了串联教材叙述而加入的解释性推导。这个边界不会削弱结论，反而让读者知道何处可以复核，何处需要在自己的实验中重新验证。如果后续课程材料更新，也可以沿着同一证据链替换来源，而不必重写整章结构。
\section{课程资料与回看路径}
教材正文已经把主要概念、公式和实现逻辑组织成连续叙述；本节只提供回到课程材料的阅读路径。读者复习时可以先完成前文练习，再按下面的时间段或资料组核对细节。
\subsection{视频回看路径}
\begin{longtable}{p{0.20\linewidth}p{0.46\linewidth}p{0.26\linewidth}}
\toprule
\textbf{时间} & \textbf{关键词} & \textbf{复习任务}\\
\midrule
00:00:05-00:06:25 & going, talk, scale, scaling, think, we're, learning & 回看定义、例子、推导或 caveat 在该段中的位置。\\
00:19:14-00:26:00 & scaling, they're, then, going, sparsity, learning, different & 回看定义、例子、推导或 caveat 在该段中的位置。\\
00:39:25-00:45:53 & think, maybe, scaling, then, scale., scale, over & 回看定义、例子、推导或 caveat 在该段中的位置。\\
00:51:54-00:57:45 & muon, think, then, it's, scale, matrix, very & 回看定义、例子、推导或 caveat 在该段中的位置。\\
01:10:16-01:16:41 & going, scaling, learning, very, it's, over, different & 回看定义、例子、推导或 caveat 在该段中的位置。\\
\bottomrule
\end{longtable}
这些时间段用于定位视频中的讲解顺序。严肃引用时，应回到视频片段确认讲者表述、板书或幻灯片内容。
回看视频时，不必逐字抄录字幕。更有效的方法是把每个时间段判定为定义、例子、推导、代码解释或 caveat，再把它放回本章对应小节。
\subsection{课程资料阅读路径}
\begin{longtable}{p{0.14\linewidth}p{0.48\linewidth}p{0.30\linewidth}}
\toprule
\textbf{资料组} & \textbf{标题/页块} & \textbf{关键词}\\
\midrule
official\_01 & Lecture 11 S C A L I N G – C A S E S T U D Y A N D D E TA I L S CS336 & scaling, scale, initialization, optimizers, cs336, motivation, best\\
official\_02 & Working through some detailed, public scaling recipes 1. MiniCPM 2. DeepSeek & scaling, minicpm, stabilize, many, working, detailed, public\\
official\_03 & Optimal LR According to muP – optimal learning rate should be (roughly) stable. Is it? & scaling, batch, learning, size, optimal, analysis, rate\\
official\_05 & Observation 3 (?) robustness Generalization to MoEs Generalization to other datasets & scaling, scale, update, parametrization, muon, gains, recall\\
official\_06 & What is muP not robust to? Exotic optimizers There are other, exotic optimizers based ... & scaling, robust, exotic, optimizers, strong, weight, decay\\
\bottomrule
\end{longtable}
课程资料通常给出更稳定的标题、公式、图或代码结构；视频则补充动机和口头解释。两者合起来构成本章的事实边界。
阅读资料时，可以把每个资料组拆成三个对象：一个定义，一个公式或伪代码动作，一个边界条件。这样比逐字翻译页面或脚本更容易形成可用知识。
\section{实践说明、历史位置与风险}
Scaling laws 进阶与实践 在 CS336 的课程结构中连接基础机制和规模化问题。基础机制解释方法为什么成立；规模化问题检验它在更大模型、更长上下文、更复杂数据或真实 serving workload 下是否仍成立。
历史位置不等于模型年表。更重要的是识别问题如何迁移：早期方法解决小规模建模问题，现代方法常常解决同一问题在大规模数据、GPU 集群、长上下文或人类偏好约束下的新形态。
\begin{itemize}
\item 资料版本限制（version risk）：课程页、公开视频和配套资料可能在学期中更新；引用时应注明访问日期或资料版本。
\item 测量口径限制（measurement risk）：FLOPs、tokens/s、benchmark score、reward 等数字只有在协议完整时可比较。
\item 规模外推限制（scaling risk）：小模型或小 batch 的机制不必然外推到 frontier-scale。
\item 数据分布限制（data risk）：数据来源、过滤和去重策略会改变模型行为，也会改变 evaluation contamination 风险。
\end{itemize}
\section{复习要点}
下面的问题把本章内容压缩为可反复检查的条目。每个条目都应能回到前文的解释、公式、伪代码或例题。
\begin{longtable}{p{0.16\linewidth}p{0.38\linewidth}p{0.38\linewidth}}
\toprule
\textbf{条目} & \textbf{对象} & \textbf{检查动作}\\
\midrule
条目 1 & forecasting（预测） & 定义它；写出一个最小例子；说明一个 failure mode；指出它影响哪个资源或评估账本。\\
条目 2 & overtraining（过训练/超额训练） & 定义它；写出一个最小例子；说明一个 failure mode；指出它影响哪个资源或评估账本。\\
条目 3 & data-constrained scaling（数据受限缩放） & 定义它；写出一个最小例子；说明一个 failure mode；指出它影响哪个资源或评估账本。\\
条目 4 & muP & 定义它；写出一个最小例子；说明一个 failure mode；指出它影响哪个资源或评估账本。\\
条目 5 & emergence（涌现） & 定义它；写出一个最小例子；说明一个 failure mode；指出它影响哪个资源或评估账本。\\
\bottomrule
\end{longtable}
\subsection{综合自测}
\begin{itemize}
\item 我能否不用课程原句，独立解释本章中心问题？
\item 我能否把本章至少一个公式改写成可运行的估算函数？
\item 我能否指出一个看似改进但实际转移成本的方法？
\item 我能否回到视频和课程资料，找到支持本章关键论断的位置？
\item 我能否设计一个最小 ablation，验证本章一个 caveat？
\end{itemize}
\section{总结与延伸}
\subsection{本章小结}
\begin{itemize}
\item Scaling laws 进阶与实践 的核心不是记忆术语，而是理解对象、变量和机制之间的关系。
\item 视频给出讲解顺序，课程资料提供可核对的公式、图、代码或页面结构。
\item 复习时应把每个术语连接到至少一个变量、一个实现动作和一个失败模式。
\end{itemize}
\subsection{延伸解释}
本章中的延伸解释用于连接课程材料中的概念，例如说明公式里的 proxy、解释系统 tradeoff，或把多个材料片段串成一个完整推理链。若需要严肃引用，应回到课程视频和配套资料核对原始表述。
\subsection{练习}
\begin{itemize}
\item 定义题：用中英双语解释本章任意五个重要术语，并说明它们属于 compute、memory、data、optimization、evaluation 还是 deployment 账本。
\item 推导题：选择本章一个公式，逐个解释符号、单位和近似假设，并说明哪个变量最难在真实系统中测量。
\item 实现题：把本章伪代码改写成一个最小 Python sanity check，不要求训练大模型，但要输出可检查的中间量。
\item 资料题：从“课程资料与回看路径”中选择一个视频时间段，回到原始视频核对讲者例子，并记录是否存在字幕误差。
\item 评估题：为本章方法设计一个 ablation 或 benchmark，说明控制变量、数据版本、指标和 failure criteria。
\item 边界题：说明本章一个结论在更长 context、更大 batch、更多 GPU 或更复杂数据分布下可能如何失效。
\item 综合题：把本章内容和前一章或后一章连接起来，写出一个端到端训练/推理 pipeline 中的依赖关系。
\item 批判题：指出一个容易被营销材料夸大的说法，并用本章的资源账本或评估协议反驳它。
\end{itemize}
\end{document}
